%% Page 3: why the name, the four properties, and the arithmetic setup (§3.1).
\crumbs{§2 Four attractive properties \,›\, §3 Arithmetic reasoning \,›\, 3.1 Experimental setup}
\insection{method,setup}

%% ---- marks on the original page
\bracket{1}{35.9}{58.6}
\bracket{2}{70.9}{82.1}
\bracket{3}{83.3}{98.4}
\bracket{4}{99.6}{110.8}
\bracket{5}{112.0}{123.2}
\bracket{6}{139.5}{171.2}
\bracket{7}{184.0}{208.6}
\bracket{8}{210.9}{226.0}
\bracket{9}{228.2}{254.8}
\highlight{101.7}{70.9}{142.3}{74.4}
\highlight{105.4}{83.3}{134.8}{86.8}
\highlight[hlpink]{103.2}{94.9}{139.0}{98.4}
\highlight[hlgreen]{72.5}{107.3}{134.8}{110.8}
\highlight{110.8}{112.0}{131.7}{115.5}
\highlight[hlpink]{155.0}{243.6}{177.8}{247.2}
\highlight[hlpink]{38.1}{247.5}{65.6}{251.0}

\block[y=36]{%
\num{1}\ \textbf{Why ``chain of thought'' and not ``solution''}: it reads like a worked solution, but it is produced
\textbf{before} the answer and imitates step-by-step thinking.}

\block[y=62]{%
\heading{Four reasons it is attractive}
\num{2}\ \textbf{Decomposition}: a multi-step problem becomes intermediate steps, so harder problems get
\textbf{more computation} (more tokens written).}

\block[y=83.3]{%
\num{3}\ \textbf{Interpretability}: we see \textbf{how} the answer was reached and can debug where it went wrong.
\begin{warning}
The paper itself says fully characterising the model's computations \hl[hlpink]{\orig{remains an open question}}:
the written chain is not guaranteed to be what the model actually computed.
\end{warning}}

\block{%
\num{4}\ \textbf{Generality}: maths, commonsense, symbols, in principle \textbf{any task people solve in words}.\quad
\num{5}\ \textbf{Cheap}: off-the-shelf large models, just add exemplars to the prompt.}

\block{%
\begin{question}
Is the extra computation of property \num{2} why it works? The ablation on page 6 says \textbf{no}: dots instead of
words do not help.
\end{question}}

\block[y=139.5]{%
\num{6}\ \textbf{§3 Arithmetic first}: simple for people, hard for language models. Claim: the 540B model + chain
of thought $\approx$ \textbf{task-specific finetuned} models, and a new state of the art on GSM8K.}

\block[y=184]{%
\num{7}\ \textbf{Five maths benchmarks}:
\begin{smalltable}{>{\bfseries}l l}
GSM8K & multi-step word problems, the hardest here \\
SVAMP & problems with varying structures \\
ASDiv & diverse problems \\
AQuA & algebraic, \textbf{multiple choice} \\
MAWPS & a mix of simpler problems \\
\end{smalltable}}

\block[y=211]{%
\num{8}\ \textbf{Baseline = standard prompting}: question/answer exemplars, the model answers directly (Figure 1, left).}

\block[y=228.2]{%
\num{9}\ \textbf{Proposed}: the same exemplars, each with a chain of thought. \textbf{Eight}, written by hand (most
datasets only have an evaluation split); full set in Appendix Table 20.
\begin{arrows}
  \item They \hl[hlpink]{\orig{did not undergo prompt engineering}}: robustness to other writers is tested in §3.4.
\end{arrows}}
